\documentclass[10pt]{article}
\usepackage[T1]{fontenc}
\usepackage[utf8]{inputenc}
\usepackage[french]{babel}
\usepackage{amsmath,amssymb,mathrsfs,mathpazo}
\usepackage{geometry}
\geometry{a4paper,margin=2cm}
\usepackage{graphicx,array,tabularx,tabularray,colortbl,multirow,multicol,booktabs,enumitem,pifont,ulem,fancybox,diagbox,cancel}
\usepackage{pstricks,pstricks-add,pst-plot,pst-tree,pst-node,pst-text,pst-func,pst-math,pst-bspline,pst-3dplot,pst-eucl,pst-solides3d}
\let\pscalculate\undefined
\usepackage{pst-bezier}
\usepackage{tikz,pgfplots,tkz-tab}
\usetikzlibrary{arrows,arrows.meta,patterns,calc,positioning,trees,decorations.markings,shapes.geometric,shapes.misc,shapes.arrows,intersections,decorations.pathmorphing,decorations.pathreplacing,backgrounds,shadows,babel}
\pgfplotsset{compat=1.18}
\usepgfplotslibrary{fillbetween}
\usepackage[np]{numprint}
\usepackage{listings,makecell,pas-tableur,twoopt}
\let\barre\undefined
\usepackage[pstricks]{tablvar}
\definecolor{darkgreen}{RGB}{0,100,0}
\providecommand{\red}{\color{red}}
\providecommand{\blue}{\color{blue}}
\providecommand{\green}{\color{green}}
\providecommand{\black}{\color{black}}
\providecommand{\white}{\color{white}}
\providecommand{\gray}{\color{gray}}
\providecommand{\cyan}{\color{cyan}}
\providecommand{\magenta}{\color{magenta}}
\providecommand{\yellow}{\color{yellow}}
\providecommand{\violet}{\color{violet}}
\providecommand{\orange}{\color{orange}}
\providecommand{\euro}{\texteuro}
\providecommand{\EUR}{\texteuro}
\providecommand{\ieme}{\textsuperscript{e}}
\providecommand{\degres}{\ensuremath{^\circ}}
\providecommand{\degree}{\ensuremath{^\circ}}
\providecommand{\vv}[1]{\overrightarrow{#1}}
\providecommand{\Coord}[1]{#1}
\providecommand{\decosix}{}
\providecommand{\starredbullet}{\ensuremath{\bullet}}
\pagestyle{empty}
\setlength{\parindent}{0pt}
\usepackage{adjustbox}\DeclareUnicodeCharacter{2236}{:}\DeclareUnicodeCharacter{2190}{\ensuremath{\leftarrow}}
\providecommand{\R}{\mathbb{R}}
\providecommand{\N}{\mathbb{N}}
\providecommand{\Z}{\mathbb{Z}}
\providecommand{\Q}{\mathbb{Q}}
\providecommand{\D}{\mathbb{D}}
\providecommand{\C}{\mathbb{C}}
\providecommand{\e}{{\mathrm{e}}}
\providecommand{\vect}[1]{\overrightarrow{#1}}
\providecommand{\vectt}[1]{\overrightarrow{\mathrm{#1}}}
\providecommand{\vvt}[1]{\overrightarrow{\mathrm{#1}}}
\providecommand{\np}[2][]{{#2~#1}}
\providecommand{\vv}[1]{{\overrightarrow{#1}}}
\renewcommand{\d}{{\mathrm{d}}}
\providecommand{\degres}{{^\circ}}
\providecommand{\degre}{{^\circ}}
\providecommand{\blue}{\color{blue}}
\providecommand{\red}{\color{red}}
\providecommand{\enskip}{\quad}
\providecommand{\llbracket}{\lbrack\!\lbrack}
\providecommand{\rrbracket}{\rbrack\!\rbrack}
\providecommand{\euro}{{\text{€}}}
\providecommand{\no}{{\text{n°}}}
\providecommand{\interval}[2]{{[#1;#2]}}
\providecommand{\intervFO}[2]{{[#1;#2[}}
\providecommand{\dotfill}{\ldots}

\providecommand{\ds}{\displaystyle}
\providecommand{\textding}[1]{#1}
\providecommand{\Oij}{\(\mathrm{O};\vec{\imath},\vec{\jmath}\)}
\providecommand{\Oijk}{\(\mathrm{O};\vec{\imath},\vec{\jmath},\vec{k}\)}
\providecommand{\Ouv}{\(\mathrm{O};\vec{u},\vec{v}\)}
\providecommand{\intervFF}[2]{\ensuremath{[#1;#2]}}
\renewcommand{\euro}{\ifmmode\text{\texteuro}\else\texteuro\fi}
\DeclareUnicodeCharacter{00B0}{\ifmmode^{\circ}\else\textdegree\fi}
\DeclareUnicodeCharacter{20AC}{\ifmmode\text{\texteuro}\else\texteuro\fi}
% Énoncés : Éducation nationale. Transcription et corrigés : APMEP, Denis Vergès et contributeurs. Découpage et classement : CoopMaths, Sébastien Lozano, Éric Elter et contributeurs. Adaptation de la présentation : Maths Explorer.
% CC BY-SA : https://coopmaths.fr/www/about/
% https://coopmaths.fr/alea/static/e3c/2020/tex/e3c_2020_00_sujet65_1_cor.tex
\begin{document}
\noindent\textit{APMEP / CoopMaths. CC BY-SA. Présentation adaptée par Maths Explorer.}\par\medskip
\subsection*{1.}

\(\textbf{Affirmation 1 :}\)

On a \( \overrightarrow{AB} \begin{pmatrix} 2 \\ 2 \end{pmatrix} \) et \( \overrightarrow{CD} \begin{pmatrix} -7 \\ 6 \end{pmatrix} \), donc :
\[
\overrightarrow{AB} \cdot \overrightarrow{CD} = 2 \times (-7) + 2 \times 6 = -14 + 12 = -2 \neq 0,
\]
le produit scalaire n'est pas nul. Les vecteurs ne sont pas orthogonaux, les droites \((AB)\) et \((CD)\) ne sont pas perpendiculaires.

\(\textbf{Affirmation 2 :}\)

Le point \(E(3 \,;\, -2)\) appartient à la droite d'équation \(y = x - 5\).

Or \( \overrightarrow{CE} \begin{pmatrix} 3 \\ 3 \end{pmatrix} \) et :
\[
\overrightarrow{CE} \cdot \overrightarrow{AB} = 2 \times 3 + 2 \times 3 = 12 \neq 0.
\]
La droite d'équation \(y = x - 5\) contient le point \(C\) mais n'est pas perpendiculaire à la droite \((AB)\).

\(\textbf{Affirmation 3 :}\)

On a :
\[
AB^2 = (4 - 2)^2 + (0 - (-2))^2 = 4 + 4 = 8.
\]
Une équation du cercle de centre \(A\) passant par \(B\) est donc :
\[
(x - 2)^2 + (y - (-2))^2 = 8 \iff (x - 2)^2 + (y + 2)^2 = 8.
\]
Affirmation vraie.

\subsection*{2.}

\(\textbf{Affirmation 4 :}\)

On dérive \(f\) comme quotient de fonctions dérivables, puisque \(x\) est non nul :
\[
f'(x) = \dfrac{\e^x \times x - 1 \times \e^x}{x^2} = \dfrac{\e^x (x - 1)}{x^2}.
\]
D'où \(f'(1) = \dfrac{\e^1 (1 - 1)}{1^2} = \dfrac{0}{1} = 0\).

\subsection*{3.}

\(\textbf{Affirmation 5 :}\)

\(\dfrac{2\pi}{5}\) radians correspondent à \(\dfrac{2 \times 180}{5} = 72\) en degré.

On est donc dans le premier quadrant : le cosinus et le sinus sont tous les deux positifs : l'affirmation est fausse.

\end{document}
